\begin{theorem}\label{thm:ok-algebraic-volume-order}
Let \(\varphi,\psi\in\PSH(X,\theta)\) have positive multiplier-ideal
volume.
Then
\[
 \lambda_n\!\left(\Delta_\nu(\theta,\varphi)\right)
     =\frac1{n!}\vol(\theta,\varphi).
\tag{55.4}\label{eq:ok-algebraic-volume}
\]
If
\(\varphi\preceq_{\cI}\psi\), meaning
\(\cI(k\varphi)\subseteq\cI(k\psi)\) for every \(k>0\), then
\[
 \Delta_\nu(\theta,\varphi)
       \subseteq\Delta_\nu(\theta,\psi).
\tag{55.5}\label{eq:ok-i-order}
\]
In particular the body depends only on the \(\cI\)-singularity type.
For a potential with minimal singularities it is the usual Okounkov body
\(\Delta_\nu(L)\).
\end{theorem}

\begin{proof}[Almost-semigroup source boundary]
The ideal inclusion gives
\(W_k(\theta,\varphi)\subseteq W_k(\theta,\psi)\) for every \(k\).
Monotonicity of the almost-semigroup body gives
\eqref{eq:ok-i-order}.

The remaining assertion is the central almost-semigroup theorem of
\cite{Xia21}, reproduced in the supplied Okounkov chapter of
\cite{XiaBook}.  We record its inputs so that the boundary is visible.
First, the multiplier-ideal family
\(\Gamma_\nu(\theta,\varphi)\) is proved to be an almost semigroup, and its
asymptotic counting volume is
\[
 \lim_{k\to\infty}k^{-n}
 h^0\!\left(X,L^k\otimes\mathcal I(k\varphi)\right)
 =\frac1{n!}\vol(\theta,\varphi).
\]
This is the multiplier-ideal asymptotic Riemann--Roch theorem.  Second, the
almost-semigroup body construction is continuous for its pseudometric and
has Euclidean volume equal to the counting volume.

For orientation, the source proves these inputs first for analytic
singularities defining a K\"ahler current.  On a log resolution
\[
                   \pi^*\theta_\varphi=[D]+R,
\tag{55.6}\label{eq:ok-log-resolution}
\]
divisorial valuations translate the ordinary Okounkov body by
\(\nu(D)\), and asymptotic Riemann--Roch computes the residual volume.  It
then passes to K\"ahler currents by quasi-equisingular approximation and
to arbitrary positive-volume potentials by
\(\varphi_\varepsilon=(1-\varepsilon)\varphi+\varepsilon\chi\).
Continuity is required both for the almost semigroups and for
\(\mathcal I\)-model volume.  Applying the two imported inputs to
\eqref{eq:ok-partial-body} gives
\eqref{eq:ok-algebraic-volume}.
\end{proof}
